\section{Persistent Topology Extension}

This section adds a persistent topological observable to the integrated state.
It is deliberately narrower than a claim that topology identifies structure:
the carrier is a fixed finite simplicial complex, and the filtration is the
only varying input in the stability result below.

\begin{definition}[Filtered finite carrier]
Let $K$ be a finite abstract simplicial complex and let $\F$ be a field. A
filtration function is a map $f:K\to\R$ satisfying $\sigma\subseteq\tau
\Rightarrow f(\sigma)\leq f(\tau)$. Its sublevel complex at $a\in\R$ is
\[
 K_a^f=\{\sigma\in K:f(\sigma)\leq a\}.
\]
The persistent signature in degree $k$ is the persistence diagram
$\Dgm_k(f)$ of the homology persistence module
$H_k(K_a^f;\F)$, and the truncated signature is
$\Pi_r(K,f)=(\Dgm_0(f),\ldots,\Dgm_d(f))$.
\end{definition}

The monotonicity condition makes every $K_a^f$ a simplicial subcomplex: if
$\tau$ is retained, every face $\sigma$ is retained as well. Thus inclusion
of sublevel complexes gives the persistence maps used in the definition.

\begin{definition}[Persistent discrepancy]
For two filtration functions $f,g$ on the same finite carrier $K$, define
\[
 \Delta_{\mathrm{pers},r}(f,g)=\max_{0\leq k\leq r}
 d_B\bigl(\Dgm_k(f),\Dgm_k(g)\bigr),
\]
where $d_B$ is the bottleneck distance. The two filtrations are compared only
after fixing the field, the carrier, the degree cutoff, and the units of the
filtration values.
\end{definition}

\begin{theorem}[Filtered stability; imported theorem]
For a fixed finite simplicial complex $K$ and monotone filtration functions
$f,g:K\to\R$,
\[
 d_B\bigl(\Dgm_k(f),\Dgm_k(g)\bigr)\leq \|f-g\|_\infty
 \quad\text{for every }k.
\]
Consequently, $\Delta_{\mathrm{pers},r}(f,g)\leq\|f-g\|_\infty$.
\end{theorem}

\begin{proof}
Put $\delta=\|f-g\|_\infty$. For every simplex $\sigma$, the inequalities
$g(\sigma)\leq f(\sigma)+\delta$ and
$f(\sigma)\leq g(\sigma)+\delta$ hold. Therefore, for every $a$,
\[
 K_a^f\subseteq K_{a+\delta}^g
 \quad\text{and}\quad
 K_a^g\subseteq K_{a+\delta}^f.
\]
The inclusions commute with the inclusions between sublevel complexes as $a$
increases. Applying homology gives a $\delta$-interleaving of the two
persistence modules in degree $k$. The algebraic stability theorem for
persistence modules gives a bottleneck distance no larger than the
interleaving distance. Hence $d_B(\Dgm_k(f),\Dgm_k(g))\leq\delta$.
Taking the maximum over $0\leq k\leq r$ proves the final inequality. The
only non-elementary step is the imported algebraic stability theorem; the
sublevel inclusions and interleaving construction are established here.
\end{proof}

\begin{proposition}[Isomorphism invariance of the persistent signature]
Let $\varphi:K\to L$ be a simplicial isomorphism and let $f:K\to\R$ and
$g:L\to\R$ satisfy $g(\varphi(\sigma))=f(\sigma)$ for every simplex
$\sigma$. Then, for every $k$,
\[
 d_B\bigl(\Dgm_k(f),\Dgm_k(g)\bigr)=0.
\]
\end{proposition}

\begin{proof}
Fix $a$. The value-preservation condition implies
$\varphi(K_a^f)=L_a^g$. Since the restriction
$\varphi_a:K_a^f\to L_a^g$ is a simplicial isomorphism, functoriality of
simplicial homology gives $H_k(K_a^f;\F)\cong H_k(L_a^g;\F)$. For
$a\leq b$, the square formed by these restrictions and the inclusion-induced
homology maps commutes. Thus the persistence modules are isomorphic and their
persistence diagrams are identical, so the bottleneck distance is zero.
\end{proof}

\begin{proposition}[The persistent signature is not complete for a labeled carrier]
Even on one fixed labeled carrier, equality of all persistence diagrams need
not determine the filtration as labeled structural data.
\end{proposition}

\begin{proof}
Let $K$ consist of two isolated vertices $x,y$ with distinct labels. Set
$f(x)=0,f(y)=1$ and $g(x)=1,g(y)=0$. Both degree-zero diagrams contain
$(0,\infty)$ and $(1,\infty)$, and every positive-degree diagram is empty.
Thus $\Pi_r(K,f)=\Pi_r(K,g)$ for every $r\geq0$. However, the only
label-preserving bijection fixes both vertices, and it does not preserve the
filtration values. Hence the persistent signature does not determine the
labeled filtered state.
\end{proof}


The extension certifies invariance and controlled sensitivity under its stated
hypotheses. It does not establish completeness, comparison of unrelated
carriers without a correspondence, or neural discovery of the carrier. Those
questions require separate assumptions or empirical validation.


\begin{proposition}[State-level persistent invariance]
Let $\Sigma$ and $\Sigma'$ be integrated states and let
$\varphi:\Sigma\to\Sigma'$ be a structural isomorphism. Then for every
degree cutoff $r$,
\[
d_B(\Dgm_k(f),\Dgm_k(g))=0\quad\text{for every }k\leq r.
\]
\end{proposition}

\begin{proof}
The filtration-preserving component of $\varphi$ satisfies
$g(\bar\varphi\sigma)=f(\sigma)$ for every simplex. Therefore the
isomorphism-invariance proposition applies to the two filtered carriers.
\end{proof}